%; whizzy paragraph -pdf xpdf -latex ./whizzypdfptex.sh
%; whizzy-paragraph "^\\\\begin{frame}\\|\\\\emtext"
% latex beamer presentation.
% platex, latex-beamer でコンパイルすることを想定。 

%     Tokyo Debian Meeting resources
%     Copyright (C) 2012 Junichi Uekawa

%     This program is free software; you can redistribute it and/or modify
%     it under the terms of the GNU General Public License as published by
%     the Free Software Foundation; either version 2 of the License, or
%     (at your option) any later version.

%     This program is distributed in the hope that it will be useful,
%     but WITHOUT ANY WARRANTY; without even the implied warreanty of
%     MERCHANTABILITY or FITNESS FOR A PARTICULAR PURPOSE.  See the
%     GNU General Public License for more details.

%     You should have received a copy of the GNU General Public License
%     along with this program; if not, write to the Free Software
%     Foundation, Inc., 51 Franklin St, Fifth Floor, Boston, MA  02110-1301 USA

\documentclass[cjk,dvipdfmx,12pt]{beamer}
\usetheme{Tokyo}
\usepackage{monthlypresentation}

%  preview (shell-command (concat "evince " (replace-regexp-in-string "tex$" "pdf"(buffer-file-name)) "&")) 
%  presentation (shell-command (concat "xpdf -fullscreen " (replace-regexp-in-string "tex$" "pdf"(buffer-file-name)) "&"))
%  presentation (shell-command (concat "evince " (replace-regexp-in-string "tex$" "pdf"(buffer-file-name)) "&"))

%http://www.naney.org/diki/dk/hyperref.html
%日本語EUC系環境の時
\AtBeginDvi{\special{pdf:tounicode EUC-UCS2}}
%シフトJIS系環境の時
%\AtBeginDvi{\special{pdf:tounicode 90ms-RKSJ-UCS2}}

\newenvironment{commandlinesmall}%
{\VerbatimEnvironment
  \begin{Sbox}\begin{minipage}{1.0\hsize}\begin{fontsize}{8}{8} \begin{BVerbatim}}%
{\end{BVerbatim}\end{fontsize}\end{minipage}\end{Sbox}
  \setlength{\fboxsep}{8pt}
% start on a new paragraph

\vspace{6pt}% skip before
\fcolorbox{dancerdarkblue}{dancerlightblue}{\TheSbox}

\vspace{6pt}% skip after
}
%end of commandlinesmall

\title{東京エリアDebian勉強会}
\subtitle{第125回 2015年4月度}
\author{野島貴英}
\date{2015年4月18日}
\logo{\includegraphics[width=8cm]{image200607/openlogo-light.eps}}

\begin{document}

\begin{frame}
\titlepage{}
\end{frame}

\begin{frame}{設営準備にご協力ください。}
会場設営よろしくおねがいします。
\end{frame}

\begin{frame}{Agenda}
 \begin{minipage}[t]{0.45\hsize}
  \begin{itemize}
   \item 注意事項
	 \begin{itemize}
	  \item 写真はセミナールーム内のみ可です。
          \item 出入りは自由でないので、もし外出したい方は、野島まで一声くださいませ。
	 \end{itemize}
   \item 事前課題発表
  \end{itemize}
 \end{minipage} 
 \begin{minipage}[t]{0.45\hsize}
  \begin{itemize}
   \item 最近あったDebian関連のイベント報告
	 \begin{itemize}
	 \item 第124回 東京エリアDebian勉強会
	 \end{itemize}
   \item Debian Trivia Quiz
   \item .debからPythonパッケージへの変遷
   \item 今後のイベント
   \item 今日の宴会場所
  \end{itemize}
 \end{minipage}
\end{frame}

\section{事前課題}
\emtext{事前課題}
{\footnotesize
\begin{prework}{ 野島 }
  \begin{enumerate}
  \item Q.hack timeに何をしますか？\\
    A. DDTSSなどなど\\
    http://ddtp.debian.net/ddtss/index.cgi/ja
  \item (オプション)Q.何について聞きたい／参加者と話をしたいですか？\\
    A. "Designed for Windows 10"で無理やり導入されるかもしれないSecure bootなPCについての傾向と対策とソフトウェア自由の確保の為の戦い。
  \end{enumerate}
\end{prework}

\begin{prework}{ まえだこうへい }
  \begin{enumerate}
  \item Q.hack timeに何をしますか？\\
    A. いくつか溜まっているパッケージのメンテナンス
  \item (オプション)Q.どこで今回の勉強会の開催を知りましたか？\\
    A. Debian JPのメーリングリスト
  \end{enumerate}
\end{prework}

\begin{prework}{ wskoka }
  \begin{enumerate}
  \item Q.hack timeに何をしますか？\\
    A. Tile-GX8036搭載機にdebianを移植しようと試みる
  \item (オプション)Q.どこで今回の勉強会の開催を知りましたか？\\
    A. その他
  \item (オプション)Q.何について聞きたい／参加者と話をしたいですか？\\
    A. 新しいアーキテクチャのdebian化についてご意見を伺いたい
  \end{enumerate}
\end{prework}

\begin{prework}{ keikurata }
  \begin{enumerate}
  \item Q.hack timeに何をしますか？\\
    A. 初心者なので見るだけでお願いします。
  \item (オプション)Q.どこで今回の勉強会の開催を知りましたか？\\
    A. その他
  \end{enumerate}
\end{prework}

\begin{prework}{ NOKUBI Takatsugu }
  \begin{enumerate}
  \item Q.hack timeに何をしますか？\\
    A. SimStringのパッケージ化、KAKASI関連
  \item (オプション)Q.どこで今回の勉強会の開催を知りましたか？\\
    A. その他のメーリングリスト
  \item (オプション)Q.何について聞きたい／参加者と話をしたいですか？\\
    A. mecab-ipadic-neologdのパッケージ化について
  \end{enumerate}
\end{prework}

\begin{prework}{ koedoyoshida }
  \begin{enumerate}
  \item Q.hack timeに何をしますか？\\
    A. DDTSSまたはPyconJP関連
  \item (オプション)Q.どこで今回の勉強会の開催を知りましたか？\\
    A. Debian JPのメーリングリスト
  \end{enumerate}
\end{prework}

\begin{prework}{ wbcchsyn }
  \begin{enumerate}
  \item Q.hack timeに何をしますか？\\
    A. 何かDebianに寄与する事。多分、翻訳かバグフィックス
  \item (オプション)Q.どこで今回の勉強会の開催を知りましたか？\\
    A. 友達や知り合いから直接
  \end{enumerate}
\end{prework}

\begin{prework}{ alohaug }
  \begin{enumerate}
  \item Q.hack timeに何をしますか？\\
    A. Gnuk 1.1.4がgpg4win/Win7で認識されないバグを追いかけます。
  \item (オプション)Q.どこで今回の勉強会の開催を知りましたか？\\
    A. Debian JPのメーリングリスト
  \end{enumerate}
\end{prework}

\begin{prework}{ henrich }
  \begin{enumerate}
  \item Q.hack timeに何をしますか？\\
    A. リリースノート訳
  \item (オプション)Q.どこで今回の勉強会の開催を知りましたか？\\
    A. Twitter(@tokyodebian)
  \item (オプション)Q.何について聞きたい／参加者と話をしたいですか？\\
    A. Jessieに対する不安
  \end{enumerate}
\end{prework}

\begin{prework}{ dictoss }
  \begin{enumerate}
  \item Q.hack timeに何をしますか？\\
    A. debian packageにおけるkernelの作法を調べる
  \item (オプション)Q.どこで今回の勉強会の開催を知りましたか？\\
    A. Debian JPのメーリングリスト
  \end{enumerate}
\end{prework}

\begin{prework}{ yy\_y\_ja\_jp }
  \begin{enumerate}
  \item Q.hack timeに何をしますか？\\
    A. DDTSS 
  \item (オプション)Q.どこで今回の勉強会の開催を知りましたか？\\
    A. その他
  \item (オプション)Q.何について聞きたい／参加者と話をしたいですか？\\
    A. DDTSSのレビューのお願い
  \end{enumerate}
\end{prework}

\begin{prework}{ Roger Shimizu }
  \begin{enumerate}
  \item Q.hack timeに何をしますか？\\
    A. jessieに関する検証・作業
  \item (オプション)Q.どこで今回の勉強会の開催を知りましたか？\\
    A. Twitter(@tokyodebian)
  \end{enumerate}
\end{prework}

}

\section{イベント報告}
\emtext{イベント報告}

\begin{frame}{第124回東京エリアDebian勉強会}

\begin{itemize}
\item 場所はスクウェア・エニックスさんのセミナルームをお借りしての開催でした。
\item 参加者は9名でした。
\item セミナ内容は岩松さんによる「Raspberry Pi 2 Model B に Debian Jessie / armhf をインストールする」でした。
\item 残りの時間でhack timeを行い、成果発表をしました。
\item 宴会の代わりに、「祥龍房 新宿イーストサイドスクエア店」で夕食会をやりました。
\end{itemize} 
  
\end{frame}

\begin{frame}{第124回東京エリアDebian勉強会(つづき)}

 セミナはDebian公式開発者の岩松さんにより、最近発売になったRaspberry Pi 2 Model BにDebian Jessieのarmhfアーキテクチャをインストールするお話でした。

\begin{itemize}
\item Raspberry Pi 2 は、Debian Jessieのarmhfアーキテクチャがそのままで動作。
\item Debianの元でUnixBenchを取ると、Raspberry Pi (armel): Raspberry Pi 2 (armhf) = 66.5:450.8となりRaspberry Piより何倍も性能が向上
\item cdebootstrapの詳しい動作・説明
\end{itemize} 

\end{frame}

\begin{frame}{第124回東京エリアDebian勉強会(つづき)}

  参加者にて、各種 Raspberry Pi が持ち込まれ、hack timeで早速インストールなど試みる方もいらっしゃいました。
\begin{center}  
\LARGE
 組み込みでもDebianを使う人がどんどん増えると良いとイイな！
\end{center}

\end{frame}
  
\section{Debian Trivia Quiz}
\emtext{Debian Trivia Quiz}
\begin{frame}{Debian Trivia Quiz}

  Debian の常識、もちろん知ってますよね?
知らないなんて恥ずかしくて、知らないとは言えないあんなことやこんなこと、
みんなで確認してみましょう。

今回の出題範囲は\url{debian-devel-announce@lists.debian.org},
\url{debian-news@lists.debian.org} に投稿された
内容などからです。

\end{frame}

\subsection{問題}

%; whizzy-master ../debianmeetingresume201311.tex
% 以上の設定をしているため、このファイルで M-x whizzytex すると、whizzytexが利用できます。
%

\santaku
{2015/2/17にて、DPLのLucas Nussbaumにより立て直しの協力募集が行われたチームは？}
{DSA team}
{partners team}
{東京エリアDebian勉強会team}
{B}
{Debian Projectには、Debian パートナーズプログラム(https://www.debian.org/partners/)という制度があります。こちらを統括しているチームとしてpartners teamがあるのですが、昨今、こちらの構成員の時間が取れない状態らしく、対応が滞っているという状況のようです。なお、DSAはThe Debian System Administratorsというチームのことでバリバリ活動されてます。東京エリアDebian勉強会teamは当勉強会幹事の頭の中の妄想上のチームですが、頑張ってるぜ！}

\santaku
{2015/2/23に、Debianがクレジットに載せられた映画があることがindenti.caのzakの投稿で判りました。何と言う名前の映画？}
{ミレニアム ドラゴン・タトゥーの女}
{Toy Story}
{Citizenfour}
{C}
{Citizenfourというエドワード・スノーデンを題材にしたドキュメンタリー映画のクレジットのSpecial ThanksとしてDebianが入りました。imdbにもcompanyとして登録されました。Citizenfourは2015年アカデミー賞のうち長編ドキュメンタリー賞を受賞しました(http://www.imdb.com/company/co0504449/, http://www.huffingtonpost.jp/tatsuhei-morozumi/citizenfour\_b\_6741756.html)ギャガ株式会社配給で将来日本でも公開されるとのことですので、将来クレジットを確認すると良いかもしれません}

\santaku
{2015/3/31にJessieのリリース目標の日がアナウンスされました。いつでしょう？}
{4/1}
{4/18}
{4/25}
{C}
{遂に4/25がJessieのリリース目標の日取りとなりました。ただ、最悪のケースとして、Jessieにリリースに支障があるような非常に重大な問題が見つかり修正が間に合わない場合は、ずれる可能性があるとは記載されています。現在uddを参照すると、キーパッケージだけでも47個ぐらいRCバグが残ったままのようですので、Fix! Fix!　なお、4/25に東京でリリースパーティーが企画されています。詳細はconnpassのDebian JPグループを見て下さいませ。}

\santaku
{2015/4/16のlucas最後のbit from DPLが流れました。こちらに記載されていたDebianからパッケージをリリースできるようにするための法的対策を検討中のソフトウェアは以下のどれ?}
{libdvdcss}
{OTR software}
{Unreal Engine4}
{A}
{なんと、libdvdcssと、ZFSが検討中らしい。もちろん、ダメになる可能性は十分にあるので、状況を静観したいと思います。libdvdcssはそのコード本体がクラッキング行為の塊なので、万一debianに入った場合、日本でdebianをミラーしても大丈夫かどうかは今後の議論ですね}

\santaku
{Debianにて開発者の属性の不均衡を正す事を目的として活動する公式チームがDebian Projectに新設したと4/16にlucasによりアナウンスがありました。なんという名前のチーム？}
{Debian Release team}
{Debian Outreach team}
{Technical Comittie}
{B}
{現在、是正すべき属性のターゲットとしては性別分布を目標としているようです。現在、Debian公式開発者の性別の割合は男性に大きく偏ってしまっています。Debianは多様性を重視するプロジェクトです。将来、性別の不均衡が是正されると良いですね。}

\santaku
{2015/4/15に選出された新DPLは誰？}
{Mehdi Dogguy}
{Gergely Nagy}
{Neil McGovern}
{C}
{Neil McGovernさんはDebian ProjectでRelease Teamで活動されるなど多数の貢献をされている方です（所信表明は：https://www.debian.org/vote/2014/platforms/neilm　）。ちなみに、2015のDebian JP Project会長は岩松さんが選出されました。おめでとうございます。}








\section{.debから Python パッケージへの変遷}
\emtext{.debから Python パッケージへの変遷}

\section{今後のイベント}
\emtext{今後のイベント}
\begin{frame}{今後のイベント}
  \begin{center}
    {\Large Debian8 "Jessie" リリースパーティ \\ ＠Tokyo \\4/25 10:30-14:00}\\
    \url{http://debianjp.connpass.com/event/14241/}
  \end{center}
\end{frame}

\begin{frame}{今後のイベント(つづき)}
\begin{itemize}
 \item 関西エリアDebian勉強会
 \item 東京エリアDebian勉強会(5/18か、5/25のどちらか)
\end{itemize}
\end{frame}

\section{今日の宴会場所}
\emtext{今日の宴会場所}
\begin{frame}{今日の宴会場所}
未定
\end{frame}

\end{document}

;;; Local Variables: ***
;;; outline-regexp: "\\([ 	]*\\\\\\(documentstyle\\|documentclass\\|emtext\\|section\\|begin{frame}\\)\\*?[ 	]*[[{]\\|[]+\\)" ***
;;; End: ***
